\section{Conclusion}
\label{sec:conclusion}

The question we began with was not which positional encoding is best, but what
property of a task predicts whether a positional encoding will help. We proposed
one such property---whether the task's positional rule is written in terms of the
sequence length---and tested it with a factorial design that separates ``failed to
fit'' from ``failed to extrapolate''. The answer was no: the coupling class did
not predict the sign of the effect, and \nope{} was never worse than an explicit
code, including on the two tasks constructed to favour its competitors. What we
can say with more confidence than before is narrower but firmer: the advantage of
omitting positional encoding survives the condition we designed to break it.

The broader methodological point is that recommendations about architectural
components are conditional objects. An aggregate ranking over a task battery is a
summary of a distribution, and if that distribution is heterogeneous the summary
can invert the sign of the effect for any given task. Reporting the conditioning
variable alongside the ranking costs little and makes the recommendation
actionable: a practitioner who knows their task's coupling class can choose a
positional encoding deliberately rather than by reputation.
